\section{Limitations and Potential Societal Impact}
\label{app:limitation}
\mind\ is designed to facilitate the development and evaluation of generalist agents for the web. Such agents hold great potential for making the web more accessible and easy to use, especially for individuals who are less familiar with information technology or have disabilities and may struggle to navigate through complex web apps and get overwhelmed by the options available.
However, there are still potential concerns and limitations regarding the current data collection, system design and safety for deployment in real world.

\noindent\textbf{Diversity and Representation in Data Collection.}
Although we strive to choose representative websites covering diverse domains, the present selection predominantly comprises English-language websites primarily used in the U.S.
Meanwhile, all our annotators are sourced through the Amazon MTurk platform, which might be biased towards a group that is more proficient in web use. 
Therefore, the tasks and websites embodied in our dataset may represent only a subset of all potential tasks that can be performed on the web.
Bearing this limitation in mind, the design of \mind\ and our data collection protocol allow for easy expansion to encompass more tasks and websites.
The inclusion of additional websites, potentially from different countries and languages, and tasks from more diverse demographics, such as individuals from different age groups, those traditionally facing web accessibility challenges, and professionals from specific domains like software development, research, law, and more, present exciting directions for future development.

\noindent\textbf{Use of Multimodal Information.}
Our current approach, \model, models the web environment using only textual context from webpage snapshots. Nevertheless, crucial information can also be gleaned from the visual representation of a rendered webpage.
While not currently utilized, we have included complete webpage snapshots in \mind, enabling rendering of the webpage for visual interpretation. The use of this multimodal information will be a viable prospect for improving model performance.

\noindent\textbf{Modeling of Interaction Dynamics.}
In \model, we encode each webpage independently at every step, with only the previous actions provided as historical context.
However, the changes of the web environment could also provide significant cues for task completion, such as the appearance of a dropdown menu following a button click.
Exploring effective ways to model such dynamic environment transformations during interaction could be an essential aspect for developing robust web agents.

\noindent\textbf{Human-Agent Interaction.}
In the current design of \mind, the user provides a single description of the task goal up front, and the agent carries out the task from start to finish.
In real-world settings, the user may wish to adjust or add task requirements in the middle, or the agent might seek user confirmation for more accurate task understanding.
Extending Mind2Web to an interactive or conversational setting, thereby allowing diverse forms of human-agent interactions, could be an interesting future direction.

\noindent\textbf{Evaluation with Offline/Online Environments.}
Following recent works~\cite{sunMETAGUIMultimodalConversational2022, Burns2022ADF}, we evaluate the system with cached offline environments, which allows us to test using snapshots of complex real-world websites.
However, a downside to this is that the task will fail immediately if an action was not cached during data collection, potentially leading to false negatives due to the existence of multiple paths for completing the same task.
As described in Appendix~\ref{app:eval_normalization}, we normalize the actions to address equivalent elements within the same page. 
In addition, we include complete network traffic in the dataset, presenting possibilities for future research to enable some degree of replay and exploration within the cached environment.
Given that \mind\ faithfully replicates real-world webpages, systems trained on the dataset should be readily transferable to live websites. 
Conducting end-to-end live evaluation on real websites with human assistance is a very promising direction that is worth exploration.

\noindent\textbf{Safety in Deployment.}
While the development of general-purpose web agents holds great potential to enhance efficiency, optimize user experiences, and promote web accessibility universally, the accompanying safety considerations for real-world deployment cannot be ignored.
These include how to effectively manage sensitive actions like financial transactions, enhancing transparency and interpretability, and keeping users in control during task execution.
Additionally, there is the risk of these agents possessing the capability to breach existing security measures such as CAPTCHA and being exploited for malicious activities, such as disseminating false information.
Therefore, it is also important for cybersecurity research to consider these potential uses and develop preemptive protective measures.



